\scanpage{15}
\deleted{On dit}\inserted{Soit} $T$ un voisinage
\deleted{tubulaire}\inserted{conique} de $Y\subset X$
dans $X$, posons
\[
(11)\qquad X^*=X\setminus\mathring T=(X\setminus T)\cup\dot T
\]
c'est une partie fermée de $X$. \deleted{On dit que}

\uline{Prop\unsure{.}} Conditions équivalentes (où $\lambda\in]0,1[$ donné)

a) $\dot T$ est une partie-bord de $X^*=X\setminus\mathring T$

b) Il existe un homéomorphisme de $X$
avec lui-$=$, transformant $T,\dot T$ en
% crops/150_scan015_P04.png: the crossed lambda-like group may absorb a parenthesis; no separate live opening parenthesis is visible.
$\deleted{\unsure{\lambda}}T_\lambda,\dot T_\lambda)$

b') Comme b, mais en supposant de plus
que l'homéomorphisme induit l'identité
sur $T_{\lambda'}$ (où \deleted{$0<\lambda'<\lambda$}
$\lambda'$ est donné, $0<\lambda'<\lambda$).

c) Il existe un voisinage \deleted{tubulaire} cônique $T'$
de $Y$ dans $X$, \deleted{tel que et \unsure{a}} tel que $T=\lambda T'$,
l'isomorphisme étant compatible avec les
structures de voisinages côniques.

On dit alors que $T$ est un vois. \deleted{tub.} conique
\uline{strict} dans $X$

\margin{Il vaut mieux de definir les voisinages côniques
\unsure{strict|stricts} comme \unsure{des} \uline{germes} de voisinage
plus grand $(1+\varepsilon)T\unsure{-R_{>0}}$
$(\varepsilon>0)$}

\uline{Corollaire} Quand $T$ est vois. \deleted{tub.} con. strict dans $X$,
$\dot T$ est la frontière de $T$ dans $X$.

\uline{Corollaire} Si \deleted{il est}
% crops/150_scan015_P11.png: additional triangular ink has an A-like outline; token/state and possible insertion-caret function remain uncertain.
\unsure{\deleted{$A$}|$A$}\inserted{$T$} voisinage
\deleted{tubulaire}\inserted{conique} de $Y$
dans $X$, les $T_\lambda$ $(\lambda\in]0,1[)$ sont des voisinages
\deleted{tubulaires}\inserted{coniques} stricts
\margin{14}
